\documentclass[../report.tex]{subfiles}
\begin{document}

Mỗi module là một circuit con trong file \file{.circ}, có tên chân cố định. Mạch SC có 11 circuit con; mạch PL
dùng lại cả 11 và thêm 7 circuit riêng: bốn thanh ghi pipeline và ba khối xử lý hazard (bảng~\ref{tab:modules}).
Chương này mô tả 11 circuit dùng chung; 7 circuit riêng của PL được mô tả ở chương~\ref{ch:pipe}, cạnh cơ chế
mà chúng phục vụ. \sig{PerfCounters} và bốn thanh ghi pipeline không có test vector riêng, chúng được kiểm tra qua
các chương trình chạy trên toàn CPU (mục~\ref{sec:vec}). Ảnh chụp lấy trực tiếp từ Logisim.

\begin{table}[!htb]
\centering\footnotesize
\caption{Các circuit con của hai CPU}\label{tab:modules}
\begin{tabular}{@{}llll@{}}
\toprule
Circuit con & Có ở & Mô tả ở mục & Test vector riêng \\
\midrule
\sig{ALU}, \sig{ALUControl} & SC, PL & \ref{sec:alu} & có \\
\sig{RegisterFile} & SC, PL & \ref{sec:regfile} & có \\
\sig{ImmGen} & SC, PL & \ref{sec:immgen} & có \\
\sig{Control} & SC, PL & \ref{sec:control} & có \\
\sig{BranchUnit} & SC, PL & \ref{sec:branch} & có \\
\sig{LoadUnit}, \sig{StoreUnit}, \sig{AddrDecode} & SC, PL & \ref{sec:mem} & có \\
\sig{IOController} & SC, PL & \ref{sec:io} & có \\
\sig{PerfCounters} & SC, PL & \ref{sec:perf} & không \\
\sig{IF\_ID}, \sig{ID\_EX}, \sig{EX\_MEM}, \sig{MEM\_WB} & PL & \ref{sec:pregs} & không \\
\sig{ForwardUnit}, \sig{IDBypass} & PL & \ref{sec:fwd} & có \\
\sig{HazardUnit} & PL & \ref{sec:loaduse}, \ref{sec:hzeq} & có \\
\bottomrule
\end{tabular}
\end{table}

\section{ALU và ALUControl}\label{sec:alu}

\sig{ALU} nhận hai toán hạng \sig{A}, \sig{B} 32 bit và mã \sig{ALUCtrl} 4 bit, cho ra \sig{Y} 32 bit. Mười phép
toán được tính song song bằng Adder, Subtractor, hai Comparator (một có dấu, một không dấu), ba cổng logic 32 bit
và ba Shifter. MUX 8:1 chọn kết quả theo \sig{ALUCtrl[2:0]}, tức funct3; hai MUX 2:1 theo \sig{ALUCtrl[3]} chọn
giữa ADD và SUB, giữa SRL và SRA. Lượng dịch chỉ lấy 5 bit thấp của \sig{B}, đúng với đặc tả RV32I. Kết quả 1 bit
của Comparator được mở rộng 0 lên 32 bit cho SLT và SLTU.

\begin{table}[H]
\centering\small
\caption{Mã hóa \sig{ALUCtrl}}\label{tab:aluctrl}
\begin{tabular}{@{}cc|cc|cc|cc|cc@{}}
\toprule
Mã & Phép & Mã & Phép & Mã & Phép & Mã & Phép & Mã & Phép \\
\midrule
0000 & ADD & 1000 & SUB & 0001 & SLL & 0010 & SLT & 0011 & SLTU \\
0100 & XOR & 0101 & SRL & 1101 & SRA & 0110 & OR & 0111 & AND \\
\bottomrule
\end{tabular}
\end{table}

\sig{ALUControl} sinh \sig{ALUCtrl} từ \sig{ALUOp}, funct3 và bit 30 của lệnh. Với lệnh R, bit 30 chọn SUB hoặc SRA.
Với lệnh I, bit 30 thuộc immediate: lệnh \texttt{addi x1, x1, -1024} có bit 30 bằng 1. Vì vậy bit 3 của
\sig{ALUCtrl} chỉ lấy bit 30 khi lệnh là R, hoặc khi lệnh là I với funct3 = 101 (SRAI). Thiếu điều kiện này thì
ADDI với immediate âm bị tính thành phép trừ; \file{t01\_alu.s} có trường hợp kiểm tra lỗi đó.

\begin{figure}[!htb]
\centering
\includegraphics[width=0.70\textwidth]{ALU_hr.png}
\caption{Datapath của circuit \sig{ALU} (\sig{op} = \sig{ALUCtrl[2:0]}, \sig{alt} = \sig{ALUCtrl[3]}; chân vào và
bộ tách \sig{ALUCtrl} ở mép trái mạch không đưa vào hình)}\label{fig:alu}
\end{figure}

\section{RegisterFile}\label{sec:regfile}

\sig{RegisterFile} có 31 thanh ghi 32 bit (x1 đến x31), hai cổng đọc tổ hợp \sig{RA1}/\sig{RD1},
\sig{RA2}/\sig{RD2} và một cổng ghi đồng bộ \sig{WA}, \sig{WD}, \sig{WE}. Mỗi cổng đọc là một MUX 32:1 có đầu vào
0 nối hằng 0, nên x0 không phải một thanh ghi thật. Cổng ghi dùng Decoder 5$\to$32 với enable bằng \sig{WE}; đầu ra
thứ $k$ nối vào chân enable của thanh ghi x$k$, nên ở cạnh lên chỉ thanh ghi được chọn chụp giá trị.

Module không tự bypass giá trị đang ghi sang cổng đọc: với \texttt{addi x1, x1, 1} trên CPU đơn chu kỳ, đường
ALU $\to$ \sig{WD} $\to$ \sig{RD1} $\to$ ALU sẽ thành vòng tổ hợp và mạch dao động. Pipeline dùng khối
\sig{IDBypass} đặt bên ngoài (mục~\ref{sec:fwd}).

\section{ImmGen}\label{sec:immgen}

\sig{ImmGen} tách lệnh thành 8 mảnh bit có vị trí cố định, ghép lại theo từng định dạng I, S, B, U, J, mở rộng dấu
lên 32 bit rồi chọn bằng MUX theo \sig{ImmSel} (0 = I, 1 = S, 2 = B, 3 = U, 4 = J). Do các mảnh immediate nằm cùng
chỗ giữa các định dạng (imm[10:5] luôn ở bit [30:25], imm[4:1] của S và B luôn ở bit [11:8]), khối này chủ yếu là
Splitter, Bit Extender và một MUX.

Offset của B và J là bội của 2 nên đặc tả không mã hóa bit 0; \sig{ImmGen} nối hằng 0 vào vị trí này. Bit 1 vẫn
được mã hóa vì đặc tả chọn đơn vị 2 byte để dùng chung với phần mở rộng C, có lệnh nén dài 2 byte
\cite{riscv-spec}. Với RV32I, đích nhảy có bit 1 bằng 1 là lệch địa chỉ lệnh; mạch không phát hiện vì
\sig{IMEM} chỉ dùng \sig{PC[11:2]} (mục~\ref{sec:limits}).

\section{Control}\label{sec:control}

\sig{Control} nhận lệnh 32 bit và sinh mọi tín hiệu điều khiển. Tín hiệu \sig{valid} = \sig{opcode[1]} AND
\sig{opcode[0]} làm enable cho Decoder 5$\to$32 chọn theo \sig{opcode[6:2]}. Mười đầu ra được dùng tương ứng mười nhóm
lệnh; mỗi tín hiệu điều khiển là OR của các nhóm cần nó, ví dụ \sig{RegWrite} = OP $\lor$ OP-IMM $\lor$ LOAD
$\lor$ JAL $\lor$ JALR $\lor$ LUI $\lor$ AUIPC (bảng~\ref{tab:ctrl}).

\begin{table}[H]
\centering\small
\caption{Tín hiệu điều khiển theo nhóm lệnh}\label{tab:ctrl}
\setlength{\tabcolsep}{4pt}
\begin{tabular}{@{}llccccccccc@{}}
\toprule
Nhóm & opcode & RW & ASel & BSel & ALUOp & Imm & MR & MW & WB & B/J/Jr \\
\midrule
OP (R) & 0110011 & 1 & 0 & 0 & 10 & -- & 0 & 0 & 0 & 0/0/0 \\
OP-IMM & 0010011 & 1 & 0 & 1 & 11 & 0 & 0 & 0 & 0 & 0/0/0 \\
LOAD & 0000011 & 1 & 0 & 1 & 00 & 0 & 1 & 0 & 1 & 0/0/0 \\
STORE & 0100011 & 0 & 0 & 1 & 00 & 1 & 0 & 1 & -- & 0/0/0 \\
BRANCH & 1100011 & 0 & 1 & 1 & 00 & 2 & 0 & 0 & -- & 1/0/0 \\
JAL & 1101111 & 1 & 1 & 1 & 00 & 4 & 0 & 0 & 2 & 0/1/0 \\
JALR & 1100111 & 1 & 0 & 1 & 00 & 0 & 0 & 0 & 2 & 0/0/1 \\
LUI & 0110111 & 1 & 2 & 1 & 00 & 3 & 0 & 0 & 0 & 0/0/0 \\
AUIPC & 0010111 & 1 & 1 & 1 & 00 & 3 & 0 & 0 & 0 & 0/0/0 \\
SYSTEM & 1110011 & \multicolumn{9}{l}{\sig{Halt} = 1, các tín hiệu khác bằng 0} \\
\bottomrule
\end{tabular}

\vspace{0.3em}
\begin{minipage}{0.95\textwidth}\footnotesize
RW = \sig{RegWrite}, Imm = \sig{ImmSel}, MR = \sig{MemRead}, MW = \sig{MemWrite}, WB = \sig{WBSel},
B/J/Jr = \sig{Branch}/\sig{Jal}/\sig{Jalr}. ASel: 0 = rs1, 1 = PC, 2 = hằng 0. WBSel: 0 = kết quả ALU,
1 = dữ liệu bộ nhớ, 2 = PC + 4.
\end{minipage}
\end{table}

Ngoài các tín hiệu trên, \sig{Control} sinh \sig{UseRs1} và \sig{UseRs2} cho biết lệnh có thật sự đọc rs1, rs2 hay
không. Pipeline dùng hai tín hiệu này để không stall khi trường rs1 của LUI hoặc JAL tình cờ trùng rd của một lệnh
load.

\section{BranchUnit}\label{sec:branch}

\sig{BranchUnit} so sánh hai toán hạng bằng một Comparator không dấu (cho phép bằng và nhỏ hơn không dấu) và một
Comparator có dấu. MUX 4:1 chọn theo \sig{funct3[2:1]}: 00 là bằng, 10 là nhỏ hơn có dấu, 11 là nhỏ hơn không
dấu. Kết quả XOR với \sig{funct3[0]}, vì RISC-V mã hóa BNE, BGE, BGEU là phủ định của BEQ, BLT, BLTU, rồi AND với
\sig{Branch} để ra \sig{Taken}.

\section{LoadUnit, StoreUnit và AddrDecode}\label{sec:mem}

\sig{AddrDecode} tách địa chỉ thành \sig{WordAddr} = \sig{Addr[13:2]}, \sig{Off} = \sig{Addr[1:0]} và
\sig{IsIO} = (\sig{Addr[31:28]} = \texttt{0xF}). \sig{RamWE} = \sig{MemWrite} AND NOT \sig{IsIO}, \sig{IoWE} =
\sig{MemWrite} AND \sig{IsIO}.

\sig{LoadUnit} dịch phải word đọc được \sig{Off}$\times$8 bit để byte cần lấy về vị trí thấp nhất, mở rộng dấu hoặc
mở rộng 0 theo 8 hay 16 bit, rồi chọn bằng MUX theo funct3 (000 LB, 001 LH, 010 LW, 100 LBU, 101 LHU).

\sig{StoreUnit} làm phép đọc, sửa, ghi. Với mặt nạ $M$ = mask $\ll$ (\sig{Off}$\times$8), trong đó mask là
\texttt{0xFF} cho SB, \texttt{0xFFFF} cho SH và \texttt{0xFFFFFFFF} cho SW, word mới là
$\mathit{New} = (\mathit{Old} \land \lnot M) \lor ((\mathit{WData} \ll \mathit{Off}\times8) \land M)$.
Ví dụ \texttt{sb} byte \texttt{0xAB} vào offset 2 của word \texttt{0x11223344} cho $M$ = \texttt{0x00FF0000} và
word mới \texttt{0x11AB3344}.

\begin{figure}[!htb]
\centering
\includegraphics[width=\textwidth]{StoreUnit_hr.png}
\caption{Circuit \sig{StoreUnit}: ghép byte hoặc halfword mới vào word cũ}\label{fig:store}
\end{figure}

\section{IOController}\label{sec:io}

\sig{IOController} hiện thực UART và GPIO theo kiểu memory-mapped (bảng~\ref{tab:memmap}). \sig{Addr[4:2]} qua
Decoder 3$\to$8 (enable = \sig{IsIO}) chọn thanh ghi I/O. Khi ghi, \sig{TxWE} = \sig{MemWrite} AND chọn 0 đẩy
ký tự \sig{WData[6:0]} ra TTY; thanh ghi \sig{GPIO\_OUT} 32 bit có enable = \sig{MemWrite} AND chọn 4. Khi đọc,
MUX 8:1 trả về ký tự từ Keyboard, thanh ghi trạng thái, \sig{GPIO\_OUT} hoặc 8 công tắc DIP. Ở mức CPU, MUX
\sig{IsIO} ? \sig{IoRData} : RAM đưa dữ liệu vào \sig{LoadUnit}, nên \texttt{lbu}, \texttt{lh} trên thanh ghi I/O vẫn đúng.

Đọc \texttt{UART\_RXDATA} lấy ký tự ra khỏi hàng đợi Keyboard, nên \sig{KbdRE} phải AND với \sig{MemRead}; nếu
không, một lệnh \texttt{add} tình cờ có kết quả \texttt{0xF0000004} cũng làm mất ký tự. Trong pipeline, việc đọc
chỉ xảy ra ở tầng MEM, sau điểm flush, nên EX/MEM mang theo trường \sig{MemRead}.

\section{PerfCounters}\label{sec:perf}

\sig{PerfCounters} có hai bộ đếm 32 bit. \sig{cycles} tăng ở mỗi cạnh lên khi CPU đang chạy (\sig{Run} = 1, tức
chưa halt). \sig{instret} tăng khi có lệnh hoàn thành (\sig{Retire} = 1): ở SC là mỗi chu kỳ, ở PL là khi lệnh ở
tầng WB có \sig{valid} = 1, nên bubble không được đếm. Bộ test đọc hai chân này để tính CPI.

\end{document}
